\begin{frame}{on-policy 라는 뜻: 행동 정책 = 목표 정책}
  \begin{block}{핵심 정리 (5.2)}
    $\varepsilon$-greedy MC 제어는 (1) 탐험을 섞은 \textbf{행동
    정책 $\pi$} 의 $Q^\pi$ 를 추정하고, (2) 거기서 $\arg\max$
    를 취해 \textbf{결정론적 정책}을 만들어낸다. ``추정
    대상''과 ``최종 배포 정책''이 살짝 달라지는 이 미묘함은
    Week 9 DQN 과 Week 11 PPO 에서도 계속 따라다닌다.
  \end{block}
  \vskip0.5em
  ``학습하는 정책($\pi$)''과 ``그 정책의 가치를 측정하는
  기준''이 \textbf{같은 정책}이라는 점이, 이번 절의 MC 제어가
  \kb{on-policy} 방법이라는 뜻. 5.3절에서
  ``\textbf{무작위로 모은 데이터로 다른 정책의 가치를} 알고
  싶다''는 \kb{off-policy} 문제로 넘어갈 때, 바로 이
  ``행동 정책과 목표 정책이 같아야 한다는'' 전제가 어떤
  불편함을 만드는지 살펴본다.
\end{frame}
